\chapter{The completion architecture}
\label{ch:completion}
\minitoc\label{the-completion-architecture}
\index{completion!architecture}\index{completing-read}\index{completion-at-point}\index{vertico}\index{corfu}\index{orderless}

% ── Opening signpost ──────────────────────────────────────────────
Chapter~\ref{ch:keybindings} showed how keystrokes reach commands;
this chapter picks up where they arrive.  Emacs has two orthogonal
completion systems that newcomers often confuse.  This chapter separates minibuffer completion
(\texttt{completing-read}) from in-buffer completion (CAPF), explains
the vertico/corfu stack, and clarifies how \texttt{orderless}
and \texttt{completion-styles} interact.

Emacs has two orthogonal completion systems that are often confused because they share some vocabulary.

\section{The two axes}\label{the-two-axes}

The two completion axes are summarised in Table~\ref{tab:completion-axes}.

\begin{sidewaystable}[p]
\centering
\caption{Minibuffer vs.\ in-buffer completion axes}
\label{tab:completion-axes}
\begin{tabularx}{\linewidth}{@{}lllX@{}}
\toprule
Axis & API & Where & Replaceable by \\
\midrule
\textbf{Minibuffer completion} & \texttt{completing-read} & Prompts, \texttt{M-x}, \texttt{C-x b} & vertico, ivy, helm, ido \\
\textbf{In-buffer completion} & \texttt{completion-at-point} (CAPF) & Inside a buffer (code) & corfu, company \\
\bottomrule
\end{tabularx}
\end{sidewaystable}

These are independent. \texttt{vertico} replaces the minibuffer UI without touching in-buffer completion. \texttt{corfu} replaces the in-buffer popup without touching the minibuffer.

\section{Minibuffer completion: the vertico stack}\label{minibuffer-completion-the-vertico-stack}

The minibuffer stack in this config:

\begin{minted}{text}
completing-read / read-from-minibuffer
    → completion-styles: orderless + basic          (HOW candidates match)
    → completion-category-overrides                 (per-category style)
    → vertico                                       (VERTICAL UI, replaces default)
    → marginalia                                    (ANNOTATIONS on each candidate)
    → consult                                       (ENHANCED COMMANDS feeding the stack)
    → embark                                        (ACTIONS on any candidate)
\end{minted}

\texttt{orderless} extends \texttt{completion-styles}: it matches space-separated patterns in any order with fuzzy/regex dispatch. It does \textbf{not} change the UI; that is vertico's job.

\texttt{consult} does not replace \texttt{completing-read}; instead, it provides richer \textbf{sources} to it. \texttt{consult-buffer} passes buffers, recent files, and bookmarks to the same \texttt{completing-read} that vertico renders.

\texttt{embark} adds a transient action layer on top of the minibuffer: while any \texttt{completing-read} is active, \texttt{C-.} opens an action menu for the current candidate. \texttt{embark-export} pushes candidates into a dedicated buffer (e.g., a \texttt{grep} buffer for text matches) where \texttt{wgrep} can make them editable in-place.

\section{In-buffer completion: the CAPF protocol}\label{in-buffer-completion-the-capf-protocol}

\texttt{completion-at-point-functions} (CAPF) is a hook whose members are functions that return completion data when point is inside a completable region. Each function either returns \texttt{nil} (not applicable here) or a list \texttt{(start end collection :props...)}.

The built-in \texttt{completion-at-point} command (\texttt{C-M-i} / \texttt{M-TAB}) iterates the hook until one function returns non-nil, then shows candidates. \texttt{corfu} intercepts at the same point and shows candidates in a popup instead.

\texttt{cape} (Completion At Point Extensions) adds extra CAPF functions: \texttt{cape-dabbrev} (dynamic abbrev from open buffers), \texttt{cape-file} (file paths), \texttt{cape-keyword} (language keywords). These are added to \texttt{completion-at-point-functions} globally or buffer-locally.

\texttt{eglot} registers its own CAPF function when it manages a buffer: that function calls the LSP server's \texttt{textDocument/completion} endpoint and returns server-provided candidates.

\subsection{Why not company?}\label{why-not-company}

Company implements its own backend protocol that is parallel to and incompatible with CAPF. corfu + cape is a lighter, CAPF-native alternative: no parallel protocol, any CAPF-aware package (eglot, elisp, tex, etc.) works without a dedicated company backend.

\section{\texttt{completion-styles} and orderless in depth}\label{completion-styles-and-orderless-in-depth}

Each element of \texttt{completion-styles} is tried in order until one matches. The built-in styles:

\begin{longtable}[]{@{}ll@{}}
\toprule
Style & Matching rule \\
\midrule
\texttt{basic} & Prefix match only \\
\texttt{partial-completion} & Partial prefix match, e.g. \texttt{/u/s/l} matches \texttt{/usr/share/lib} \\
\texttt{flex} & Fuzzy subsequence match \\
\texttt{orderless} & Space-separated, each component matched independently \\
\bottomrule
\end{longtable}

This config's setting:

\begin{minted}{elisp}
(setq completion-styles '(orderless basic)
      completion-category-overrides
        '((file (styles basic partial-completion))))
\end{minted}

File completion uses \texttt{basic} + \texttt{partial-completion} (no orderless) because path components are ordered. \texttt{partial-completion} handles \texttt{$\sim$/D/P/s} → \texttt{$\sim$/Documents/Projects/src} expansion correctly, while \texttt{orderless} would match path segments in arbitrary order, producing surprising results for file name completion.

% ── Mental model summary ──────────────────────────────────────────
\section{Mental model}
Emacs completion divides into two independent axes: minibuffer completion, where \texttt{completing-read} feeds candidates through vertico, and in-buffer completion, where CAPF functions feed candidates through corfu.  The minibuffer stack layers several cooperating packages---\texttt{orderless} handles flexible matching, vertico provides the vertical UI, marginalia adds annotations, consult supplies rich candidate sources, and embark offers contextual actions on any candidate.  The \texttt{completion-styles} list is tried in order; \texttt{orderless} covers most scenarios, but file-name completion specifically needs \texttt{partial-completion} to expand path prefixes correctly.

On the in-buffer side, CAPF functions return \texttt{(start end collection)} or \texttt{nil}, and corfu renders their output as a popup.  Corfu together with cape is fully CAPF-native, unlike company which uses a parallel, incompatible backend protocol.  Eglot (Chapter~\ref{ch:eglot-lsp}) registers its own CAPF backend so that LSP completions flow into this architecture transparently.  Meanwhile, \texttt{embark-export} can push any set of candidates into a dedicated buffer where \texttt{wgrep} makes them editable, bridging the gap between selection and bulk editing.
